\documentclass{../note}

\title{课程时间安排 HUFE}
\author{isomoes}

\begin{document}

\maketitle

\section{考纲内容}

\subsection{考核目标}
\begin{enumerate}
	\item 考核学生对各种数据结构的基本概念与基本原理的理解和掌握，以及运用数据结构知识分析问题和解决问题的能力。
	\item 考核学生对数据库系统的基本概念与基本原理的理解和掌握，以及运用数据库设计方法分析问题和解决问题的能力。
\end{enumerate}

\subsection{考核内容}
\subsubsection{数据结构}
\begin{enumerate}
	\item \textbf{绪论}\\
	      数据、数据元素、数据结构、数据类型、抽象数据类型的概念，数据的逻辑结构和存储结构，算法、算法描述和算法分析的概念。

	\item \textbf{线性表}\\
	      线性表的定义及其抽象数据类型描述，顺序表的逻辑结构定义及其基本运算，链表的逻辑结构及其基本操作。

	\item \textbf{栈和队列}\\
	      栈的结构特性、基本操作及在顺序存储结构和链式存储结构上基本运算的实现，队列的结构特性、基本操作及在顺序存储结构和链式存储结构上基本运算的实现，栈和队列的基本应用。

	\item \textbf{数组和广义表}\\
	      数组的基本概念和存储结构，广义表的定义和存储结构。

	\item \textbf{树和二叉树}\\
	      树的基本概念，二叉树的概念、性质和存储结构，二叉树的遍历，线索二叉树，哈夫曼树。

	\item \textbf{图}\\
	      图的基本概念，图的存储结构（邻接矩阵、邻接表、十字链表和邻接多重表），图的遍历，生成树和最小生成树，最短路径。

	\item \textbf{查找}\\
	      查找的基本概念，线性表的查找，二叉排序树，哈希表的查找。

	\item \textbf{内排序}\\
	      排序的基本概念，各种排序（插入排序、交换排序、选择排序、归并排序和基数排序）的基本思想和算法分析。
\end{enumerate}

\subsubsection{数据库原理}
\begin{enumerate}
	\item \textbf{绪论}\\
	      数据库的4个基本概念，数据管理技术的产生和发展，数据建模、概念模型和数据模型的三要素，数据库系统的三级模式结构，数据库的两级映像与数据独立性，数据库系统的组成。

	\item \textbf{关系模型}\\
	      关系模型的数据结构及形式化定义，关系操作，关系完整性，关系代数（传统的集合运算、专门的关系运算）。

	\item \textbf{关系数据库标准语言SQL}\\
	      数据定义、数据查询、数据更新、空值处理、视图。

	\item \textbf{数据库安全性}\\
	      数据库安全性概述，数据库安全性控制。

	\item \textbf{数据库完整性}\\
	      数据库完整性概述，实体完整性，参照完整性，用户定义完整性，完整性约束命名子句。

	\item \textbf{关系数据理论}\\
	      关系数据库规范化理论的基本概念，函数依赖的定义和函数依赖的公理系统，第一/二/三范式和BC范式，关系模式的分解。

	\item \textbf{数据库设计}\\
	      数据库设计的基本步骤及各阶段的主要任务，E-R模型及用E-R模型进行概念结构设计，逻辑结构设计。

	\item \textbf{数据库恢复和并发控制}\\
	      事务的基本概念，故障的种类，恢复的实现技术，恢复策略及具有检查点的恢复技术；并发控制的基本概念。
\end{enumerate}

\section{课程时间安排}

基于考纲内容、章节难度以及总课时压缩至15小时的实际情况，我们将课程安排为5次课，每次3小时。

\subsection{总体时间分配}
\begin{itemize}
	\item 数据结构：8.5小时（重点保留高频和高难内容）
	\item 数据库原理：6.5小时（突出关系模型、SQL和设计基础）
	\item 授课方式：5次课，每次3小时
	\item 总计：15小时
\end{itemize}

\subsection{分次授课安排（5次课）}
\begin{tabular}{|c|c|p{10cm}|}
	\hline
	\textbf{课次} & \textbf{时间} & \textbf{主要内容} \\
	\hline
	第1次课 & 3小时 & 数据结构绪论、线性表、栈和队列，完成基础概念梳理与典型题入门 \\
	\hline
	第2次课 & 3小时 & 数据结构中的数组和广义表、树和二叉树，重点讲解二叉树性质、遍历和哈夫曼树 \\
	\hline
	第3次课 & 3小时 & 数据结构中的图、查找、内排序，并进行数据结构部分阶段总结 \\
	\hline
	第4次课 & 3小时 & 数据库绪论、关系模型、SQL语言，重点突破关系代数与SQL查询 \\
	\hline
	第5次课 & 3小时 & 数据库安全性、完整性、关系数据理论、数据库设计、恢复与并发控制，并进行综合串讲 \\
	\hline
\end{tabular}

\subsection{详细时间安排}

\subsubsection{数据结构部分 (8.5小时)}

\begin{tabular}{|l|c|p{8cm}|}
	\hline
	\textbf{章节内容} & \textbf{时间分配} & \textbf{教学要点}             \\
	\hline
	1. 绪论         & 0.5小时         &
	数据结构的基本概念、逻辑结构与存储结构，算法分析入门                         \\
	\hline
	2. 线性表        & 1小时           &
	顺序表和链表的实现与基本操作，重点掌握单链表及典型题型                         \\
	\hline
	3. 栈和队列       & 1小时           &
	栈和队列的结构特点与基本操作，结合表达式求值等经典应用                         \\
	\hline
	4. 数组和广义表     & 0.5小时         &
	了解多维数组存储思想与广义表基本概念，以识记为主                             \\
	\hline
	5. 树和二叉树      & 2小时           &
	二叉树性质、遍历算法、线索二叉树与哈夫曼树，突出遍历与应用题                     \\
	\hline
	6. 图          & 1.5小时         &
	图的存储、遍历、最小生成树和最短路径，重点掌握核心算法思想                       \\
	\hline
	7. 查找         & 1小时           &
	顺序查找、折半查找、二叉排序树和哈希表，强调常见考点                           \\
	\hline
	8. 内排序        & 1小时           &
	常见排序算法的基本思想及时间复杂度比较，突出插入、交换、选择和归并排序               \\
	\hline
\end{tabular}

\subsubsection{数据库原理部分 (6.5小时)}

\begin{tabular}{|l|c|p{8cm}|}
	\hline
	\textbf{章节内容} & \textbf{时间分配} & \textbf{教学要点} \\
	\hline
	1. 绪论         & 0.5小时         &
	数据库基本概念，三级模式结构与两级映像                           \\
	\hline
	2. 关系模型       & 1小时           &
	关系模型核心概念与关系代数基本操作，重点掌握选择、投影、连接                 \\
	\hline
	3. SQL语言      & 2小时           &
	DDL、DML、DCL及重点查询语句，强化连接查询、子查询与视图                     \\
	\hline
	4. 数据库安全性     & 0.5小时         &
	数据库安全性基本概念与访问控制方法                           \\
	\hline
	5. 数据库完整性     & 0.5小时         &
	实体完整性、参照完整性和用户定义完整性的基本要求                         \\
	\hline
	6. 关系数据理论     & 1小时           &
	函数依赖、1NF/2NF/3NF和BCNF，突出范式判断与模式分解                     \\
	\hline
	7. 数据库设计      & 0.5小时         &
	E-R模型设计及概念模型向逻辑模型转换，掌握设计流程                       \\
	\hline
	8. 数据库恢复和并发控制 & 0.5小时         &
	事务ACID属性、并发控制与恢复技术的基本概念                         \\
	\hline
\end{tabular}

\end{document}
